\chapter{Conclusion and Future Work}\label{ch:conclusion}

Section~\ref{sec:summary_contributions} retraces the preceding chapters; Section~\ref{sec:limitations} records the limitations of the theory and evidence; and Section~\ref{sec:future_directions} collects the directions this work opens.

\section{Summary of Contributions}\label{sec:summary_contributions}

The starting point was the tension Chapter~\ref{ch:introduction} made precise: monotonicity and convexity make a model's reasoning legible, but imposing them at training time restricts expressivity. The resolution is not to constrain the network but to \emph{decompose} it --- to exhibit structure a trained \ac{relu} network already contains, without altering a single prediction.

Chapter~\ref{ch:decompositions} developed this into a general theory. A \emph{valid decomposition} replaces each layer map by a pair of coordinated split-streams whose difference reproduces the network exactly, and the chapter populated the definition with instances of increasing structure. The mixing family $\fmix$ interpolates between trivial and fully separated splits; its endpoint $\eta = 1$ recovers the classical \ac{dc} decomposition of a \ac{relu} network into a difference $f = g - h$ of two convex, monotone split-streams, routing the weight parts $\Wsig{l}{+}, \Wsig{l}{-}$ into separate channels \cite{hiddenMono}. Beyond the piecewise-linear case, Theorem~\ref{thm:smooth_dc_split} established a \ac{dc}--monotone splitting for every activation with integrable curvature --- GELU \cite{Hendrycks2016}, Swish \cite{Ramachandran2017_swish}, Mish, tanh, sigmoid, softplus --- with a canonicity statement singling out the split of minimal curvature, and a bounded-variation extension recovering \ac{relu} as the Dirac limit. The stopping decomposition $\fstop$ added an early-termination channel, and the regularized family $\freg$ showed that a convex regularizer on the stop/continue decision generates smooth decompositions systematically, the softplus decomposition $\fsoft$ arising as the binary-entropy subcase --- placing the softplus replacement of \textcite{Dombrowski2020} inside the framework rather than beside it. The chapter closed with the stabilization that makes split passes usable at depth and a full \ac{cnn} extension: convolutions, gated pooling, batch normalization, and residual connections.

Chapter~\ref{ch:game_theory} changed the reading of these objects without changing them. Building on the \ac{relu}-network games of \textcite{Gaubert2025}, the stopping decomposition was reinterpreted as a two-player zero-sum \ac{sg} on a doubled state graph: weight signs route between the players $P^+$ and $P^-$, gradients and $\alpha\beta$-\ac{lrp} attributions emerge as discounted occupation measures under equilibrium policies, and the deformation parameters of Chapter~\ref{ch:decompositions} --- temperature, risk aversion, stabilizers --- become principled modifications of the game. Chapter~\ref{ch:trajectory} takes the decisive step: the equilibrium policies induce a genuine \emph{probability distribution} $\Pi$ over complete backward trajectories, and this distribution --- not any pixel-space projection --- is the backward computation. The comparison with the independent \ac{nrm} framework of \textcite{Xiong2026nrm} sharpened the stakes: both factorize the backward computation as a measure over paths, but \ac{nrm} keeps one state per neuron at the price of a \emph{signed} measure, whereas the doubled graph pays a factor of two in state space for an honest probability measure. That trade makes the full statistical toolbox applicable: the Hellinger distance between two backward computations is computable in one linear-time pass, and the parameter-randomization sanity check of \textcite{Adebayo2018sanity} lifts from attribution maps to trajectory space, where Theorem~\ref{thm:rand_stopping} shows Bhattacharyya mass decays geometrically in the number of randomized layers --- detecting collapse even when heatmaps remain deceptively similar \cite{sixt2020whenexplanationslie}.

Chapter~\ref{ch:saliency} converted structure into explanations. SplitCAM, SplitLRP, and SplitGrad apply established attribution principles to the decomposed streams, separating factual (positive) from counterfactual (negative) evidence at every layer; the difference-of-convex input neural networks (DCINNs) of Section~\ref{sec:dcinns} added a self-explainable counterpart trained directly in split form. The tropical foundations (Section~\ref{sec:tropical_polyhedra_foundations}) supply the geometric picture: through the correspondence between \ac{relu} networks and tropical rational maps \cite{Zhang2018}, the split-streams $g$ and $h$ acquire Newton polytopes, and the \emph{imaginative polytope} of Section~\ref{sec:imaginative_polytopes} interprets what the decomposition adds to the geometry. Chapter~\ref{ch:transformers} developed how gated splits of matrix and Hadamard products extend the framework to multiplicative architectures \cite{Vaswani2017}, attention entering the game view as an oracle routing policy realized on ViT-B/16 \cite{Dosovitskiy2021vit}.

Chapter~\ref{ch:experiments} assembled the empirical evidence. The split methods were evaluated on VGG16 and ResNet18 on ImageNet-S; the revised game study covers VGG16, ResNet50, and ViT-B/16 on ImageNet-S, with a ViT-B/16 transfer study on Pascal VOC. No single method dominates every metric. On ViT-B/16, one adapted $\alpha\beta$-\ac{lrp} configuration outperforms all evaluated baselines in attribution localization, and this lead transfers to Pascal VOC without re-tuning. The \ac{lrp} baselines remain stronger on faithfulness, while the anchor configurations are particularly robust. Independent validation selections for localization and faithfulness, bootstrap confidence intervals, and Holm-corrected paired tests make these trade-offs explicit. Hellinger diagnostics complement pixel-level comparisons: they distinguish trained from randomized calculations even where heatmaps remain similar, increase with Gaussian input noise, and separate intra- from inter-class pairs for at least one game on each tested architecture.

Taken together, the thesis substantiates its central claim: monotone, convex, game-theoretic, and polyhedral structure need not be trained into a network --- it is hidden in every trained \ac{relu} network, and a single decomposition framework makes all four readings available at once.

\section{Limitations}\label{sec:limitations}

The results above come with qualifications that should be stated plainly.

\paragraph{Computational overhead.} A split pass propagates two coordinated streams per channel and, currently, materializes the split network with duplicated weight matrices. Memory and compute grow by a small constant factor --- harmless for the \acp{cnn} here, but an obstacle for very large models. The overhead is an implementation artefact, not a mathematical necessity: the split can be realized through masked views of the original weights, though this thesis has not delivered such an implementation.

\paragraph{Stabilization hyperparameters.} The stabilization of Section~\ref{sec:stabilization_dc} is essential at depth and introduces choices --- thresholds, shifting schemes, and in the game view the temperatures $\tau$, $\theta$ and the risk-aversion parameter $\lambda$ --- that the theory constrains but does not determine. The revised experiments use declared validation grids and independent localization and faithfulness selections. Their differing winners show that there is no universally optimal setting; transfer to Pascal VOC is demonstrated for configurations fixed on ImageNet-S. The same holds for layer selection in SplitCAM and SplitLRP: different layers excel at different metrics.

\paragraph{Rigor of the transformer extension.} Chapter~\ref{ch:transformers} was deliberately framed as outlook. The attention representation is conditional on the realized attention pattern, LayerNorm statistics, and GELU gates. It describes relevance routing through Value paths with a detached QK-softmax reference; it does not recover the full unconditioned derivative through Q and K. These choices give a representation of a specified backward rule and do not establish that the architecture determines a unique explanatory game. The ViT-B/16 results are therefore evidence for the \emph{usefulness} of the extension, not its uniqueness or optimality.

\paragraph{Benchmark scope.} The empirical evidence is specific. All numbers are inherited verbatim from the two underlying papers, so architectures (VGG16, ResNet18, ResNet50, ViT-B/16), datasets (ImageNet-S, PASCAL VOC), and metric selections are those of the original studies; nothing was re-run. Localization against human masks is only one axis, and sharper attributions can discard network-relevant evidence outside the annotated object. All experiments concern vision models; whether the gains transfer to other modalities or into downstream task gains is untested. Finally, the trajectory-level diagnostics were validated on sanity checks and controlled perturbations --- dedicated calculation-level benchmarks do not yet exist.

\section{Future Research Directions}\label{sec:future_directions}

Four directions appear most promising.

\paragraph{A signed-measure calculus on the doubled graph.} Chapter~\ref{ch:trajectory} left the synthesis with the \ac{nrm} framework \cite{Xiong2026nrm} as an open invitation. The two generalize in orthogonal directions: \ac{nrm} contributes a marginalization calculus attributing arbitrary sets of neurons across layers by sum-product message passing, the game view probabilistic semantics, deformation parameters, and a metric on backward computations. Running \ac{nrm}'s relevance laws \emph{on the doubled state graph}, where the measure is a probability law, would give joint and conditional relevances genuine probabilistic meaning, let higher-order relevances be compared across models with the Hellinger machinery, and replace the normalization axiom with the emergent per-row discount $\gamma^{(l)}_j$ of the \ac{sg}. Such a synthesis would also supply the calculation-level benchmarks Section~\ref{sec:limitations} noted are missing.

\paragraph{Regularized decompositions for smooth activations.} Theorem~\ref{thm:smooth_dc_split} guarantees canonical \ac{dc}--monotone splits for GELU, Swish, and their relatives, and the regularized family $\freg$ of Section~\ref{sec:regularized_decompositions} explains softplus as an entropy-regularized gate --- but the two are not yet joined. The natural conjecture is that each smooth activation of the theorem arises from a convex regularizer $\Omega$ on the stop/continue decision, extending the entire pipeline --- stopping game, trajectory distribution, Hellinger diagnostics, SplitLRP --- verbatim to the GELU- and Swish-based architectures that dominate modern pretraining. Since transformers already use GELU, this is the shortest route to putting the transformer extension of Chapter~\ref{ch:transformers} on the footing of Chapter~\ref{ch:decompositions}.

\paragraph{DCINNs at scale.} The MNIST experiments of Section~\ref{sec:dcinns} showed that networks trained directly in split form separate factual from counterfactual attributes with striking clarity, but the evidence stops at small models. Scaling DCINNs raises optimization and theory questions: whether the difference-of-monotone parametrization trains stably at ImageNet scale, whether \ac{dc}-programming techniques \cite{LeThi2018DCA} or ideas from \acp{icnn} \cite{Amos2017} accelerate training, and how the accuracy--explainability trade-off behaves with capacity. A positive answer would upgrade the framework from a post-hoc lens into a design principle for self-explainable models.

\paragraph{Transformers, \acsp{llm}, and efficient split passes.} Three engineering-flavoured directions complete the picture. Extending the multiplicative splits of Chapter~\ref{ch:transformers} from encoder-style \acp{vit} to autoregressive \acp{llm} requires handling causal attention and the key--value cache inside the split pass, making trajectory-level diagnostics available where explanation quality is hardest to assess. Implementing split passes through masked views rather than materialized duplicates would remove the memory overhead of Section~\ref{sec:limitations}, and the nonnegativity of the split weights invites hardware-level exploitation. The observed layer dependence of attribution quality calls for principled layer selection or multi-layer aggregation, for which the trajectory distribution offers the right vocabulary: layers are intermediate marginals of one measure $\Pi$.

\bigskip

The unifying lesson is that a trained network is a richer object than the function it computes. One set of weights carries a difference of convex monotone parts, a zero-sum game between cooperating adversaries, a probability distribution over backward trajectories, and a pair of Newton polytopes --- none of which requires changing the network. Explaining a model, on this view, is not approximation after the fact but an act of \emph{decomposition}: choosing which of the structures already present one wishes to see.
